%%%%%%%%%%%%%%%%%%%%%%%%%%%%%%%%%%%%%%%%%%%%%%%%%%%%%%%%%%%%%%%%%%%%%%%%%%%%%%%%%%%%
% PREÁMBULO MAESTRO PARA "EL CÓDIGO DE LA REALIDAD"
% WASAFF CONSULTING
% Estilo: Profesional, minimalista, escala de grises
%%%%%%%%%%%%%%%%%%%%%%%%%%%%%%%%%%%%%%%%%%%%%%%%%%%%%%%%%%%%%%%%%%%%%%%%%%%%%%%%%%%%

% ==================================================================================
% 1. CODIFICACIÓN E IDIOMA
% ==================================================================================
\usepackage[T1]{fontenc}
\usepackage[utf8]{inputenc}
\usepackage[spanish,es-lcroman,es-tabla]{babel}
\deactivatetilden

% ==================================================================================
% 2. TIPOGRAFÍA
% ==================================================================================
\usepackage{lmodern}
\usepackage{microtype}
\usepackage{ellipsis}
\usepackage{setspace}
\setstretch{1.15}

% Configuración de fuentes KOMA-Script
\setkomafont{disposition}{\normalcolor\bfseries}
\addtokomafont{chapter}{\Huge}
\addtokomafont{section}{\Large}
\addtokomafont{subsection}{\large}
\addtokomafont{subsubsection}{\normalsize}

% ==================================================================================
% 3. GEOMETRÍA DE PÁGINA
% ==================================================================================
\usepackage[
    top=2.8cm,
    bottom=2.5cm,
    inner=3.0cm,
    outer=2.2cm,
    headheight=14pt,
    headsep=0.8cm,
    footskip=1.2cm
]{geometry}

% ==================================================================================
% 4. PALETA DE COLORES (escala de grises)
% ==================================================================================
\usepackage{xcolor}
\definecolor{wasaff-black}{gray}{0.10}
\definecolor{wasaff-darkgray}{gray}{0.25}
\definecolor{wasaff-midgray}{gray}{0.50}
\definecolor{wasaff-lightgray}{gray}{0.85}
\definecolor{wasaff-silver}{gray}{0.95}
\definecolor{wasaff-accent}{gray}{0.35}

% ==================================================================================
% 5. MATEMÁTICAS
% ==================================================================================
\usepackage{amsmath}
\usepackage{amssymb}
\usepackage{amsthm}
\usepackage{mathtools}
\usepackage{bm}

% Teoremas y entornos
\theoremstyle{definition}
\newtheorem{definicion}{Definición}[chapter]
\newtheorem{teorema}{Teorema}[chapter]
\newtheorem{lema}{Lema}[chapter]
\newtheorem{proposicion}{Proposición}[chapter]
\theoremstyle{remark}
\newtheorem{observacion}{Observación}[chapter]

% ==================================================================================
% 6. GRÁFICOS, TABLAS Y FLOTANTES
% ==================================================================================
\usepackage{graphicx}
\usepackage{booktabs}
\usepackage{array}
\usepackage{multirow}
\usepackage{tabularx}
\usepackage{float}

% ==================================================================================
% 7. LISTADOS DE CÓDIGO (listings)
% ==================================================================================
\usepackage{listings}
\renewcommand{\lstlistingname}{Listado de Código}
\renewcommand{\lstlistlistingname}{Índice de Listados de Código}

\lstset{
    language=Python,
    basicstyle=\ttfamily\small,
    keywordstyle=\bfseries\color{wasaff-black},
    commentstyle=\itshape\color{wasaff-midgray},
    stringstyle=\color{wasaff-darkgray},
    backgroundcolor=\color{wasaff-silver},
    frame=single,
    rulecolor=\color{wasaff-lightgray},
    framerule=0.4pt,
    numbers=left,
    numberstyle=\tiny\color{wasaff-midgray},
    numbersep=8pt,
    tabsize=4,
    showstringspaces=false,
    breaklines=true,
    breakatwhitespace=true,
    literate=
        {á}{{\'a}}1 {é}{{\'e}}1 {í}{{\'i}}1 {ó}{{\'o}}1 {ú}{{\'u}}1
        {Á}{{\'A}}1 {É}{{\'E}}1 {Í}{{\'I}}1 {Ó}{{\'O}}1 {Ú}{{\'U}}1
        {ñ}{{\~n}}1 {Ñ}{{\~N}}1 {ü}{{\"u}}1 {Ü}{{\"U}}1
        {¿}{{?`}}1 {¡}{{!`}}1
}

\lstnewenvironment{codigo}[1][]{%
    \lstset{#1}%
}{}

% ==================================================================================
% 8. RECUADROS (tcolorbox)
% ==================================================================================
\usepackage[skins,breakable]{tcolorbox}

\newtcolorbox{notatecnica}{
    colback=wasaff-silver,
    colframe=wasaff-lightgray,
    fonttitle=\bfseries\sffamily\small,
    title={Nota Técnica},
    attach boxed title to top left={yshift=-2mm,xshift=5mm},
    boxed title style={colback=wasaff-silver,colframe=wasaff-lightgray},
    breakable
}

\newtcolorbox{industrial}{
    colback=white,
    colframe=wasaff-midgray,
    fonttitle=\bfseries\color{white},
    coltitle=white,
    colbacktitle=wasaff-midgray,
    title={Contexto Industrial},
    attach boxed title to top left={yshift=-2mm,xshift=5mm},
    boxed title style={colback=wasaff-midgray,colframe=wasaff-midgray},
    breakable
}

% ==================================================================================
% 9. MINTED (compatibilidad con capítulos existentes)
% ==================================================================================
\usepackage{minted}
\usepackage{newfloat}
\DeclareFloatingEnvironment[
    fileext=lol,
    listname={Índice de Listados de Código},
    name=Listado de Código,
    within=chapter
]{codelisting}

% ==================================================================================
% 10. BIBLIOGRAFÍA
% ==================================================================================
\usepackage[
    backend=biber,
    style=numeric-comp,
    sorting=nyt,
    maxbibnames=6,
    minbibnames=3,
    maxcitenames=2,
    giveninits=true,
    terseinits=true,
    doi=true,
    url=false,
    eprint=false,
    isbn=false
]{biblatex}
\addbibresource{Bibliografia/BibliografiaMaestra.bib}

% ==================================================================================
% 11. HIPERVÍNCULOS
% ==================================================================================
\usepackage[
    colorlinks=true,
    linkcolor=wasaff-darkgray,
    citecolor=wasaff-darkgray,
    urlcolor=wasaff-midgray,
    bookmarksopen=true,
    bookmarksopenlevel=1,
    pdftitle={El Código de la Realidad: Simulación Multifísica desde los Primeros Principios para la Ingeniería Industrial},
    pdfauthor={Felipe Wasaff},
    pdfsubject={Simulación Computacional por Elementos Finitos para Ingeniería Industrial},
    pdfkeywords={FEM, FEniCS, Multifísica, Simulación, Elementos Finitos, Ingeniería Industrial},
    pdfcreator={LuaLaTeX + KOMA-Script},
    pdfproducer={Wasaff Consulting}
]{hyperref}

% ==================================================================================
% 12. REFERENCIAS INTELIGENTES
% ==================================================================================
\usepackage[spanish]{cleveref}
\crefname{lstlisting}{listado}{listados}
\Crefname{lstlisting}{Listado}{Listados}
\crefname{codelisting}{listado}{listados}
\Crefname{codelisting}{Listado}{Listados}
\crefname{equation}{ecuación}{ecuaciones}
\Crefname{equation}{Ecuación}{Ecuaciones}
\crefname{figure}{figura}{figuras}
\Crefname{figure}{Figura}{Figuras}
\crefname{table}{tabla}{tablas}
\Crefname{table}{Tabla}{Tablas}
\crefname{chapter}{capítulo}{capítulos}
\Crefname{chapter}{Capítulo}{Capítulos}
\crefname{section}{sección}{secciones}
\Crefname{section}{Sección}{Secciones}

% ==================================================================================
% 13. ENCABEZADOS Y PIES DE PÁGINA
% ==================================================================================
\usepackage[automark]{scrlayer-scrpage}
\pagestyle{scrheadings}
\clearpairofpagestyles
\lehead{\leftmark}
\rohead{\rightmark}
\lefoot{\pagemark}
\rofoot{\pagemark}
\renewcommand*{\chapterpagestyle}{plain}

% ==================================================================================
% 14. COMANDOS PERSONALIZADOS
% ==================================================================================
\newcommand{\fenics}{\textsc{FEniCS}}

% Conjuntos numéricos
\newcommand{\R}{\mathbb{R}}
\newcommand{\N}{\mathbb{N}}

% Operadores diferenciales
\newcommand{\diff}{\mathop{}\!\mathrm{d}}
\newcommand{\grad}{\nabla}
\newcommand{\divg}{\nabla\cdot}
\newcommand{\laplacian}{\nabla^2}

% Vectores y normas
\newcommand{\vc}[1]{\bm{#1}}
\newcommand{\norm}[1]{\left\lVert#1\right\rVert}
\newcommand{\abs}[1]{\left\lvert#1\right\rvert}
\newcommand{\inner}[2]{\left\langle#1,#2\right\rangle}

% Subíndices y símbolos
\newcommand{\ped}[1]{_{\text{#1}}}
\newcommand{\celsius}{^{\circ}\text{C}}

% Espacios de Sobolev
\newcommand{\Huno}{H^1(\Omega)}
\newcommand{\Hunocero}{H^1_0(\Omega)}
\newcommand{\Ldos}{L^2(\Omega)}

% ==================================================================================
% 15. UTILIDADES FINALES
% ==================================================================================
\usepackage{csquotes}
\usepackage{lipsum}
